
\documentclass{article}
\usepackage{colortbl}
\usepackage{makecell}
\usepackage{multirow}
\usepackage{supertabular}

\begin{document}

\newcounter{utterance}

\centering \large Interaction Transcript for game `cladder', experiment `full\_v1.5\_default', episode 3221 with Qwen3.5{-}9B.
\vspace{24pt}

{ \footnotesize  \setcounter{utterance}{1}
\setlength{\tabcolsep}{0pt}
\begin{supertabular}{c@{$\;$}|p{.15\linewidth}@{}p{.15\linewidth}p{.15\linewidth}p{.15\linewidth}p{.15\linewidth}p{.15\linewidth}}
    \# & \multicolumn{2}{c}{Player} && \multicolumn{2}{c}{Game Master} \\
    \hline

    \theutterance \stepcounter{utterance}  
    & & & \multicolumn{4}{p{0.6\linewidth}}{
        \cellcolor[rgb]{0.9,0.9,0.9}{
            \makecell[{{p{\linewidth}}}]{
                \texttt{\tiny{[P1$\langle$GM]}}
                \texttt{You are an expert in causal inference. The following question is not a typical commonsense query, but rather a meticulously designed question created by a professor specializing in causal inference, intended to assess the students' mastery of the course content.} \\
\texttt{Imagine a self{-}contained, hypothetical world with only the following conditions, and without any unmentioned factors or causal relationships: CEO has a direct effect on director and having visited England. Having visited England has a direct effect on employee. Director has a direct effect on employee. For people who have not visited England and directors who don't sign termination letters, the probability of employee being fired is 9\%. For people who have not visited England and directors who sign termination letters, the probability of employee being fired is 54\%. For people who have visited England and directors who don't sign termination letters, the probability of employee being fired is 46\%. For people who have visited England and directors who sign termination letters, the probability of employee being fired is 90\%. For people who have not visited England, the probability of director signing the termination letter is 17\%. For people who have visited England, the probability of director signing the termination letter is 27\%. For people who have visited England, would it be less likely to see employee being fired if the person had not visited England?} \\
\texttt{Start your answer with 'Yes' or 'No', followed by additional reasoning or evidence to support your explanation.} \\
            }
        }
    }
     \\ \\

    \theutterance \stepcounter{utterance}  
    & \multicolumn{4}{p{0.6\linewidth}}{
        \cellcolor[rgb]{0.561,0.737,0.561}{
            \makecell[{{p{\linewidth}}}]{
                \texttt{\tiny{[P1$\rangle$GM]}}
                \texttt{No} \\
\\ 
\texttt{**Reasoning:**} \\
\\ 
\texttt{To answer this question, we must determine the causal effect of the variable "Visited England" on the outcome "Employee Fired" for the specific population of people who **have visited England**. This requires calculating the counterfactual probability: What would be the probability of an employee being fired if a person who *has* visited England had instead *not* visited England?} \\
\\ 
\texttt{Let us define the variables:} \\
\texttt{*   \$V\$: Visited England (\$V=1\$ if yes, \$V=0\$ if no).} \\
\texttt{*   \$D\$: Director signs termination letter (\$D=1\$ if yes, \$D=0\$ if no).} \\
\texttt{*   \$Y\$: Employee fired (\$Y=1\$ if yes, \$Y=0\$ if no).} \\
\\ 
\texttt{The causal structure is: \$V \textbackslash{}rightarrow D \textbackslash{}rightarrow Y\$.} \\
\texttt{The probability of \$Y\$ depends on \$V\$ only through \$D\$. Therefore, the probability of being fired given \$V\$ and \$D\$ is denoted as \$P(Y|V, D)\$.} \\
\\ 
\texttt{**Step 1: Identify the observed probabilities for the target population (\$V=1\$).**} \\
\texttt{From the prompt:} \\
\texttt{*   \$P(Y=1 | V=1, D=0) = 46\textbackslash{}\%\$} \\
\texttt{*   \$P(Y=1 | V=1, D=1) = 90\textbackslash{}\%\$} \\
\texttt{*   \$P} \\
            }
        }
    }
    & & \\ \\

    \theutterance \stepcounter{utterance}  
    & & & \multicolumn{2}{p{0.3\linewidth}}{
        \cellcolor[rgb]{0.9,0.9,0.9}{
            \makecell[{{p{\linewidth}}}]{
                \texttt{\tiny{[GM$|$GM]}}
                \texttt{No} \\
            }
        }
    }
    & & \\ \\

    \theutterance \stepcounter{utterance}  
    & & & \multicolumn{2}{p{0.3\linewidth}}{
        \cellcolor[rgb]{0.9,0.9,0.9}{
            \makecell[{{p{\linewidth}}}]{
                \texttt{\tiny{[GM$|$GM]}}
                \texttt{yes} \\
            }
        }
    }
    & & \\ \\

    \theutterance \stepcounter{utterance}  
    & & & \multicolumn{2}{p{0.3\linewidth}}{
        \cellcolor[rgb]{0.9,0.9,0.9}{
            \makecell[{{p{\linewidth}}}]{
                \texttt{\tiny{[GM$|$GM]}}
                \texttt{game\_result = LOSE} \\
            }
        }
    }
    & & \\ \\

\end{supertabular}
}

\end{document}
